\section{The Propers: Detailed Commentary}
\subsection{A prayer before the dispute}
The Entrance takes its biblical basis from Psalm 16 in Vulgate numbering, Psalm 17 in Hebrew numbering. In the full Psalm the speaker appeals for a just hearing, asks that his steps remain firm and begs protection from enemies. Its closing hope is to behold God's glory. The Entrance selects the appeal for hearing and the image of shelter; it does not rehearse the whole case for vindication. The worshipper therefore enters as someone who needs God, before hearing of either a victorious king or a disputed tax.

Augustine assigns the prayer to Christ together with the Church, his body. His explanation of verse 6 notices an unusual order: the speaker cries because God has heard. Augustine understands present confidence as the fruit of earlier help, when the prayer was weaker. Prayer is not a force that an autonomous speaker brings to bear on a reluctant God. The ability to pray strongly already bears the mark of mercy. That is a concrete reason why dependence and confidence need not oppose each other.\footnote{Augustine, \work{Expositions on the Psalms}, XVI, §§1,6,8,14; NPNF I.8, English heading ``Psalm XVII.''}

The pupil of the eye suggests something small, precious and carefully protected. Augustine gives it a Christological reading through Christ's humanity; he reads the sheltering wings through merciful love. Aquinas, commenting directly on this Psalm, begins from the care involved in protection and offers several spiritual applications of the wings, including charity and the cross. These applications deepen the image rather than identifying an animal anatomy in God. Protection under the cross cannot mean that Christ's members never suffer. It gives them a place from which to endure suffering without making their own strength their refuge.\footnote{Aquinas, \work{Super Psalmos}, Psalm 16, §§3--4, Corpus Thomisticum, paragraph 86901 and following division.}

\subsection{The will that the Collect brings before God}
The Collect, identified by its Latin incipit in the text section, concerns devoted willing and sincere service of divine majesty. Its position matters: before the assembly hears a command about what to render, it asks for the disposition that can render rightly. The question is not only whether a duty can be named. A divided person may know the right action yet perform it for vanity, fear or advantage.

The petition thus gives an interior dimension to the Gospel's public dilemma. Sincere service cannot be manufactured by flattering the one addressed, as Jesus' questioners attempt. Nor does the prayer let the assembly treat good intentions as a private possession already secured. Asking God for readiness makes room for the change that the readings will demand. This is an account of the prayer's subject and liturgical function; its approved English is not printed or replaced by a new translation.

\subsection{Cyrus: named, strengthened, and still unknowing}
Isaiah 44 sets the prophecy against manufactured idols. Wood that becomes an object of worship cannot become the Creator who acts and calls by name. Near that chapter's close Cyrus is associated with Jerusalem's restoration. Isaiah 45 then addresses the Persian king as the Lord's anointed, speaks of the strength given him, and places that strength at the service of Jacob and Israel. The appointed selection retains both the honour and the unsettling qualification: Cyrus does not know the one who equips him.

The title raises a real exegetical difficulty. If ``anointed'' points to Christ, should the passage be read directly as an address to Jesus? Jerome refuses that conclusion. He first insists on the proper name Cyrus, against readings that turn it into ``Lord.'' He then explains the royal designation by the insignia of kingship. Most decisively, the assertion that the addressee has not known God cannot simply be transferred to the Son through whom the world is reconciled. The negative clause controls the identification just as much as the title does.\footnote{Jerome, \work{Commentary on Isaiah}, XII, Isaiah 45:1--7, DCO Latin \work{Pars 3}, pp.126--129, especially 127--128. This electronic typesetting is not a checked historical facsimile.}

Jerome nevertheless calls Cyrus a type and image of the true Christ. A type is a person or event whose real, limited action anticipates a fuller divine work. It need not share every characteristic of the fulfilment. Cyrus's liberating office makes the comparison intelligible; his ignorance marks its limit. The result is stronger than either a flat political reading or a reading that erases the king. God works through this particular ruler, and the ruler is not thereby made the final saviour.

Aquinas likewise distinguishes the king's dignity from a claim that he was actually anointed with sacred oil. In his exposition of Isaiah 45, the benefits of divine assistance have more than one intended fruit: liberation for the people, recognition of God, and the knowledge that reaches beyond Israel. He says that Cyrus remained in idolatry. The usefulness of an instrument and the conversion of its user cannot therefore be treated as the same event. A successful task is no proof that the person accomplishing it has rightly understood his life.\footnote{Aquinas, \work{Expositio super Isaiam ad litteram}, chapter 45, Corpus Thomisticum, paragraph 85005.}

The selected text moves outward. God acts for Jacob and Israel so that knowledge of the one Lord reaches from sunrise to sunset. Israel's particular election and the world's knowledge do not cancel each other. The Lord's care for a people is the means by which his identity becomes known more widely. Nothing in the passage makes the Persian king's religion a second path of equal divine authority. The same oracle that gives him a role denies the existence of another God.

\subsection{All peoples give the Creator glory}
Psalm 96 in Hebrew numbering, 95 in the Vulgate, answers the first reading with universal praise. Its selected verses call for a new song, proclaim divine glory among the nations, distinguish the Creator from idols, summon the families of peoples to worship and announce equitable government. The response from verse 7b concentrates this outward movement in a repeated act of giving God glory. The Lord who can employ a foreign ruler is not one tribal deity competing with other deities.

Augustine reads the new song through new love and the building of a house throughout the world. His account joins praise and mission. To tell the nations of God's glory is also to refuse the preacher's temptation to make himself the subject of the proclamation. Augustine's inference starts with the object of the verb: whose glory is being declared? A splendid proclamation can become false in its direction when its servant receives the honour meant for God. The Psalm therefore tests the motives of religious speakers as well as the idols of those addressed.\footnote{Augustine, \work{Expositions}, XCV, §§1--6,9,11,14--15; NPNF I.8, English heading ``Psalm XCVI.''}

His Christological reading also joins the Creator's greatness to Christ's humility. In §11 his received text includes ``from the wood,'' which he interprets of the cross. That phrase is not part of this Sunday's appointed Psalm. The strophe also omits verse 10b, the clause concerning the world's stability; Augustine's discussion of that clause belongs to the full Psalm. The selected words proclaim divine kingship and equitable judgment. Christian reception then asks how that kingship is shown in Christ.

John Paul II makes that movement explicit in his audience on this Psalm. Sanctuary worship opens into the Lord's universal rule, and prayer purifies faith from idols. He identifies the cross interpretation as a Christian rereading and ends by invoking Mark's account of the Son of Man who serves. His conclusion relates kingship to self-gift; it does not turn the ancient song into a literal narrative of crucifixion. This becomes particularly fruitful if the Mark Communion option is used.\footnote{John Paul II, General Audience, 18 September 2002, §§1--5, especially §§4--5.}

\subsection{A community whose faith has weight}
The opening of 1 Thessalonians names Paul, Silvanus and Timothy and addresses a church ``in God'' the Father and the Lord Jesus Christ. Thanksgiving follows before instruction. Paul remembers the work of faith, the labour of love and the endurance of hope. These are not three flattering labels placed over an otherwise ordinary association. Each virtue has a recognizable movement: faith acts, love spends itself, and hope lasts.

Chrysostom explains why thanksgiving to God and praise of the people belong together. Their conduct gives Paul something for which to thank God; thanksgiving teaches the hearers not to claim the whole achievement for themselves. Yet Chrysostom does not dissolve their conduct into an appearance. Faith's work is seen in steadfastness under pressure; love involves actual labour; hope is tested by prolonged affliction. The surrounding chapter describes the reception of the word amid suffering, an explanation of the reading's virtues even though that verse lies beyond the appointed endpoint.\footnote{Chrysostom, \work{Homilies on 1 Thessalonians}, I, on 1:1--7, NPNF I.13. The Sunday reading ends at 1:5b, before the clause describing the missionaries' conduct.}

Election enters this thanksgiving as something Paul recognizes through effects. Aquinas explicitly denies that God's choice has been earned by antecedent human merit. He then explains the evidence in both hearer and preacher. He allows different distributions of the references to power, Spirit and fullness: signs and spiritual efficacy, or the preacher's virtuous life, the Spirit and complete instruction. These alternatives keep the passage from becoming a slogan that identifies divine power with one spectacular manifestation alone.\footnote{Aquinas, \work{Super I Thessalonicenses}, chapter 1, lectio 1, Corpus Thomisticum, paragraph 87862.}

The form of address also matters beside the Gospel. Caesar's image identifies a coin; being ``in God'' identifies this community. Its distinction is relational before it is an achievement or a social rank. Paul does not praise the Thessalonians for commanding the city. He thanks God for a life capable of costly fidelity within it. The continuation in 1:9--10 speaks of turning from idols and waiting for the risen Son; that larger horizon explains how hope can endure without being measured by immediate success. It remains context, not an addition to the reading.

\subsection{A light held from another}
The acclamation selects only Philippians 2:15d and 16a. Its image of light and its reference to the word of life stand within Paul's exhortation to obedience after the account of Christ's humility and exaltation. A star does not need to become the whole sky in order to be visible. Likewise a Christian community can shine within the world without identifying its own influence with God's kingdom.

Aquinas distinguishes light by essence from light by participation. God is light in himself; the saints receive something of the light of the Word. Holding Christ's word therefore explains rather than competes with their brightness. The moral consequence is that visibility cannot be detached from its source. A community may become highly visible through controversy or prestige while obscuring the word that gives life.\footnote{Aquinas, \work{Super Philippenses}, chapter 2, lectio 4, paragraph 87839.}

Chrysostom's homily on the wider passage connects light with a life made blameless and with willing service instead of murmuring. His explanation of ``word of life'' stresses a life-giving principle or seed. The NPNF editor prefers an interpretation referring to the Gospel word and notes the difference. The editor's comment is not a second statement by Chrysostom. Aquinas explicitly identifies Christ's word; Chrysostom's accent here falls upon the life visible in the saints. Both aid the connection between receiving and shining, but they reach it by distinguishable explanations.\footnote{Chrysostom, \work{Homilies on Philippians}, VIII, on 2:12--16, NPNF I.13, including the editor's note at ``word of life.''}

\subsection{A trap answered through its own evidence}
Matthew places the tribute question among the Jerusalem controversies following Jesus' entry and temple teaching. The narrated alliance consists of the Pharisees' disciples and the Herodians. Its flattery is exact enough to sound like confession: Jesus teaches truthfully and does not defer to appearances. Yet the speakers attempt to use that truthfulness to force a dangerous public answer. Their attitude shows why accurate religious language and sincere service are different things.

Chrysostom describes the double trap. An answer against tribute can be made into an accusation before Roman power; an answer in its favour can discredit Jesus before the people. Jesus first exposes the intention, then requests the money. The questioners identify its image and inscription themselves. Their answer becomes evidence against the terms in which they posed the dilemma. Jesus has not merely escaped a clever conversation. He has required people who tried to control him to acknowledge a fact already present in their common life.\footnote{Chrysostom, \work{Homilies on Matthew}, LXX, §§1--2, NPNF I.10.}

The verb ``render'' is decisive for Chrysostom: the answer concerns repayment of what is due, rather than an arbitrary gift. At the same time, he limits the civil demand to what does not injure godliness. An impious requirement does not become righteous because Caesar demands it. This double affirmation rules out two easy evasions. Religious allegiance cannot excuse every refusal of civic duty; civic office cannot turn every demand into a moral obligation.

Aquinas first explains the literal civil exchange through the use of the coin. He then expressly turns to a mystical reading: the soul bears God's image and must be rendered to God. The coin is an artefact marked by earthly authority; the person is not made by Caesar. This inference draws the whole self into the second half of Jesus' answer. It is not a report that Jesus held up a second object, nor does it erase the money in his hand. The literal duty remains while the spiritual claim reaches the person who performs it.\footnote{Aquinas, \work{Super Matthaeum}, chapter 22, lectio 2, Corpus Thomisticum, paragraph 87389.}

\subsection{The offering and the servant's purification}
The Prayer over the Offerings relates the gifts through which worship is offered to the purification of those serving the mysteries. A person can possess something to offer and still need cleansing. The act of giving is therefore not a certificate of innocence. In the Gospel, the possession of a coin exposes a question; at the altar, the possession of gifts becomes an occasion to receive help.

The liturgical movement is especially searching for those who serve visibly. The act itself is good, yet the servant's motives and habits remain open to healing. This is neither contempt for human work nor the claim that an imperfect person must cease every good act. It places work in a relation of dependence: service is offered through gifts received, and service itself is accompanied by a plea for purification. The prayer's wording and source limits remain distinct from this theological explanation.

\subsection{Communion through the Psalm of mercy}
The Psalm option draws on Psalm 32:18--19 in Vulgate numbering, 33:18--19 in Hebrew numbering. The full hymn contrasts God's counsel with the insufficiency of kings, armies and horses. Its movement toward the Lord's merciful gaze thus answers a concrete temptation: confusing the instruments of protection with the power that can finally save. A great army may be an instrument within history; its greatness cannot abolish mortality.

Augustine understands the famine in the verse through deprivation of the Word and everlasting truth. Those who presume upon their own strength discover that even their strength fails when cut off from righteousness. His exposition then turns to patient waiting for incorruptible nourishment. The passage is not a promise that believers will always have abundant provisions. It exposes a hunger that provisions alone cannot satisfy.\footnote{Augustine, \work{Expositions}, XXXII, §§16--22; NPNF I.8, English heading ``Psalm XXXIII.''}

Aquinas treats mercy's effects as both rescue and strengthening. He distinguishes bodily rescue, rescue from sin and final rescue, and explains nourishment bodily, spiritually and sacramentally. His explicit sacramental reference gives the Communion setting a precise theological connection. He also insists that fear and hope belong together: fear without hope tends to despair, hope without fear to presumption. Mercy is neither a permission to disregard God nor a verdict that frailty makes return impossible.\footnote{Aquinas, \work{Super Psalmos}, Psalm 32, §§15--17, especially §16, paragraph 87103.}

\subsection{Communion through the Son's ransom}
The Mark option selects the ransom clause from 10:45. The preceding service clause and the dispute over rank belong to its Gospel context, not to the antiphon's complete appointed wording. In that context Jesus is going toward Jerusalem; James and John seek places of honour, and the others become indignant. The Lord answers the competition by distinguishing the domination associated with rulers from the service his disciples must learn. His own self-gift supplies the measure.

Bede's direct exposition of Mark makes humility, rather than power, the path to the summit of virtue. Christ gives his own action as an example for disciples who might disregard words. Bede then says that the Son of Man accepted the form of a servant to shed his blood for the world. In the same paragraph he explains ``for many'' as those who were willing to believe. The worldwide reach of the gift and its believing reception stand together in his account.\footnote{Bede, \work{In Marci evangelium expositio}, III, Mark 10:42--45, PL 92 (Paris, 1862), cols.236A--B, especially 236B.}

That conjunction prevents two reductions. The Son's ransom is not merely a lesson in helpfulness: he gives the saving gift. Nor can the gift's universal horizon be made to mean that every person already answers it in faith. Bede names the believer's response without making faith a price paid in advance for redemption. His brief gloss does not by itself explain the origin of faith's willing; the account of grace and willing in Aquinas's Philippians exposition addresses that further question.

Chrysostom, commenting on the parallel in Matthew 20:28, draws the moral argument from the distance Christ descends. The Lord's death for enemies overturns the apostles' competition for status. The one who could claim the highest place takes the servant's task. This argument is sharper than a general preference for courtesy: the saving act changes the disciples' measure of greatness. Chrysostom here comments on Matthew's parallel, while Bede comments directly on Mark.\footnote{Chrysostom, \work{Homilies on Matthew}, LXV, §§4--5, NPNF I.10.}

``Ransom'' also raises the question of release and payment. Gregory Nazianzen asks to whom Christ's blood could be offered as a ransom. He rejects the devil's entitlement to receive God and refuses an account in which the Father delights in or demands blood. He speaks instead of the Father's acceptance, the Incarnation's sanctifying work and liberation through the Son. Aquinas uses the language of superabundant satisfaction and a price paid to God, while denying that humanity ever ceased to be subject to God's power. What was lost was union in charity.\footnote{Gregory Nazianzen, \work{Oration} XLV, §§21--24, especially §22, NPNF II.7; Aquinas, \work{Summa theologiae} III, q.48, a.4, especially the body and replies, and a.5 opening.}

These are different positive accounts. Their shared refusal of a devil entitled to payment does not make Gregory's question and Aquinas's answer interchangeable. Gregory prevents the image of ransom from implying a divine appetite for violence. Aquinas explains how the voluntary loving gift can restore the relation broken by sin. Both make Christ's deed more than an example that human effort could reproduce on its own. The Communion excerpt names that saving gift before the worshipper can turn it into a programme of imitation.

\subsection{Present help and an eternal end}
The Prayer after Communion places present participation on the way toward eternal gifts. Reception is therefore both real and unfinished. The assembly has not received an empty sign, but neither has its pilgrimage ended. Its hope reaches beyond the next political success, the next practical achievement or the next feeling of religious assurance.

This ending joins the two questions with which the readings began. Earthly goods can serve a true purpose without becoming ultimate goods. A human response can be truly active without becoming its own source. Present help matters because the goal is greater than the present; preparation matters because the gift calls a person onward. The concluding prayer sends the assembly into duties that remain real, with a hope those duties cannot exhaust.
